\documentclass[11pt,a4paper]{article}
\usepackage[margin=1in]{geometry}
\usepackage[T1]{fontenc}
\usepackage{hyperref}

\title{PNP Canonical Report Status Notice}
\author{PNP public-review successor boundary}
\date{2026-06-27}

\begin{document}
\maketitle

\section*{Coordinate}
\texttt{PNP-HISTORICAL-REPORT-SANITIZED-2026-06-27-01}

\section*{Current status}
The repository contains a machine-checkable proof-certificate stack for a claimed P = NP route.
This root canonical report source now records the current public-review successor boundary.
Public theorem emission is disabled.

\begin{verbatim}
publicTheoremEmissionAllowed = false
finalTheoremReady = false
activeFinalNodeIds = []
remainingBlockers = [
  "Release.UnrestrictedFinalSoundness",
  "ExternalReview.Acceptance"
]
\end{verbatim}

\section*{Supersession of the historical report}
Earlier sealed-report artifacts and earlier revisions of this root report contained direct theorem-emission prose.
Those historical artifacts are represented by the repository supersession audit and are not the current active public-release boundary.
This document supersedes the previous root report surface as a status notice.

\section*{Verification}
The current machine-verification entry point is:

\begin{verbatim}
npm run pnp:verify
\end{verbatim}

The expected verifier output path is:

\begin{verbatim}
artifacts/pnp-verify-all/latest-verdict.json
\end{verbatim}

\section*{Non-claims}
This notice does not activate public theorem emission.
This notice does not set \texttt{finalTheoremReady}.
This notice does not clear \texttt{Release.UnrestrictedFinalSoundness} or \texttt{ExternalReview.Acceptance}.
This notice does not treat historical theorem-emission prose as the current public theorem.

\section*{Machine surfaces}
The current successor stack binds the report status through:

\begin{verbatim}
PNP_STATUS.json
report-bindings/HISTORICAL_REPORT_SUPERSESSION.json
pcc-historical-report-supersession0.mjs
scripts/pnp-verify-all.mjs
\end{verbatim}

\end{document}
